\documentclass{article}
\usepackage{amsmath,amssymb}
\title{A canonical semantic target from the generated state completion}
\date{15 November 2026}
\begin{document}\maketitle
\paragraph{Status.} Scoped construction with independent written-proof review. Its purpose is to test nonvacuity of the sufficient higher evaluation doctrine using the general generated completion itself, not a concrete implementation. It does not assert an unrestricted state endofunctor on arbitrary tricategories or a free-algebra classification.

\paragraph{Previously constructed input.} Start with the original bicategorical process input $K_0$, state bicategory $S$, pure actions and directed parallel hypotheses of the generated-tower construction. Retain all the chosen normalization data. The normalized process tricategories are $T_n$, and each next raw state construction has homwise biadjoint equivalences
\[
 N_n:\mathcal R_n(A,B)\simeq T_{n+1}(A,B).
\]
Here $\mathcal R_n$ consists of actual outer-state bodies in $T_n$, all their native squares, and their complete process three-cell witnesses. The chronological held-body assembly on these raw homs has already been identified coherently with transport of the normalized assembly. For every monotone index map $\theta:n\to m$, the specified $R_\theta:T_n\to T_m$ preserves the given pure and directed parallel structures, with the independently established indexing comparisons, modifications, perturbations and equations. No image of a raw witness is inferred solely from its outer state shadow.

\paragraph{The candidate target and its local evaluation.} Take $D=T_1$. This is a particular general categorical target, not an arbitrary process tricategory. Let $\mu:2\to1$ be the unique map. Define the local target uniformity evaluator by
\[
 A_D=R_\mu N_1:\mathcal R_1(A,B)\longrightarrow T_1(A,B).
\]
The domain includes every target body in $T_1$ at an outer state, not merely the image of one chosen seed evaluator. Thus no smaller unexplained body class is used. $N_1$ transports the full raw witness; $R_\mu$ subsequently joins the two retained state coordinates. Local functoriality and pseudofunctor constraints are compositions of those already constructed. The fixed-state operation $B_X$ is the restriction along the actual identity-state inclusion $C_X$, not a separately chosen object function.

The raw target held sequential and parallel operations are the explicit chronological operations on $\mathcal R_1$. Their action on full witnesses and their coherent identification under $N_1$ are part of the generated-tower assembly results. Composing these identifications with the sequential and parallel comparisons of $R_\mu$ provides $A_D$'s assembly comparisons. Their associativity, unit and four-process equations are the corresponding equations of those two constructions. Directed interchange is never inverted. Consequently the target doctrine's assembly is obtained from independently constructed source operations, rather than assuming the sought evaluator as data.

\paragraph{Unit recovery.} A body $f$ of $T_1$ has an inner state coordinate $Z$ and a $K_0$ body. Closing it at outer unit state gives the retained two-coordinate body $(I,Z)$, before joining. The merge gives the total state $Z\otimes I$, with its actual action frames. The admitted right unitor, including the process body square, supplies $v_f:B_I(f)\simeq f$. Its three-cell naturality uses the full native witness for a map between such $f$. The identity and assembly unit equations follow from the source indexing-unit comparisons and their stated coherence. This concerns the fixed-state inclusion at $I$; it does not quotient arbitrary state endomorphisms $I\to I$ to identities.

\paragraph{Joining from the source diagram.} A body $f$ at two outer states $X,Y$ in $T_1$ has a third retained inner coordinate $Z$. Evaluating $B_XB_Y(f)$ first merges the inner $Z,Y$ block and then joins with $X$. Evaluating $B_{Y\otimes X}(\bar f)$ first joins the two outer coordinates and then joins with $Z$. The two results are represented by the two parenthesizations of $Z\otimes Y\otimes X$, but equality of these labels is not the proof. Use the actual composition comparisons of the two index-map paths $3\to2\to1$, with normalization frames, to supply the pseudonatural joining equivalence $j$.

At a general two-boundary native square, the machine component of that indexing comparison is its actual transported process square. Its native-arrow and three-cell naturality supply exactly the target two-boundary joining comparison and its full-witness equation. This is why keeping the source comparison as a tritransformation, rather than just a state associator, is essential. Explicitly, an arbitrary target square $\alpha$ also carries an inner state map $a:Z\to Z'$, besides $k:Y\to Y'$ and $h:X\to X'$. The left route retains $(a\otimes k)\otimes h$; the joined route retains $a\otimes(k\otimes h)$. The component of $j$ is the actual associator $a_{Z,Y,X}$, with its admitted process square. Its pseudonaturality at $(a,k,h)$, pasted with the original $K_0$ square in chronological order, supplies the required target three-cell. A full witness retains all three state two-cell components and the original $K_0$ equation. Thus this verifies arbitrary inner maps, not only the identity-inner-state case.

Compatibility with independent sequential and parallel assembly follows from the already constructed structured indexing comparisons and their four-process equations. On three outer closure states, the source indexing modification gives the target triple-joining three-cell. Its pentagon and unit equations are the source higher-index coherence equations after the same normalization. More concretely, the triple-closure comparison is the state pentagonator on the four objects $Z,U,Y,X$, including the inner state $Z$. The four-closure equation is the monoidal-bicategory pentagonator coherence on five objects: $Z$ and the four outer states. The unit instances use the corresponding state unit modifications. All are supplied with their admitted $K_0$ action squares. This direct identification does not erase any source indexing perturbation's interface components. These are equations in $T_1$ obtained only after native process admission; no faithfulness assertion for an arbitrary semantic target is used.

\paragraph{What finite induction evaluates.} Let $\iota:K_0\to T_1$ be the unit-state pure inclusion, and run the conditional finite-stage construction with this seed and the preceding canonical target data. Let $!:n\to1$ denote the unique monotone map, including the empty map at $n=0$. The candidate comparison is
\[
 F_n\simeq R_!:T_n\to T_1,
\]
where $T_0\simeq K_0$ uses the fixed unit-action equivalence. At depth zero it is the actual seed unit comparison. At the next step, the inductive comparison is transported through the raw native-square construction, retaining its full witness, then through $N_1$ and $R_\mu$. The resulting merge first groups the already accumulated inner block and then the new outer state. The indexing comparison identifies this with the single total merge; the inserted seed-unit coordinate is removed with the specified admitted unit comparison. Native naturality and assembly compatibility are inherited from those same index comparisons. Hence the identification is not deduced from equality of total-state labels alone.

At $n=1$ this comparison recovers the identity of $T_1$ up to the chosen unit-state frames. Here the seed inserts an inner unit, so the comparison removes $I\otimes X$ by the left unitor. This differs from fixed-outer-unit recovery $B_I(f)$, which removes $Z\otimes I$ by the right unitor. At general $n$ it is not asserted to be an equivalence of process tricategories: joining distinct retained coordinates can forget their factorization. The notation $\simeq$ compares the two functors, not the source and target tricategories.

\paragraph{Audit boundary.} The full two-boundary rectangle and joining coherence have been independently checked using the explicit inner-state calculation above. This is a written proof, not mechanized certification. The construction proves a canonical evaluation on the generated completion, not existence in every target, evaluator uniqueness, arbitrary shared-state synchronization, or instantaneous trace.
\end{document}
